\section*{अध्याय \ref{ch:g7:how-animals-breathe} --- जंतु साँस कैसे लेते हैं}

\begin{solution}{exo:g7:how-animals-breathe:1}
एक बड़ी, पतली, नम विनिमय-सतह, जहाँ \omterm{prop:g4:journey-of-food:absorb}{रक्त} (या शरीर का तरल) आस-पास से
मिलता है। बड़ी: ताकि कुल मिलाकर पर्याप्त पार हो; पतली: ताकि पार करना
आसान हो; नम: क्योंकि गैसें घुली हुई ही पार करती हैं।
\end{solution}

\begin{solution}{exo:g7:how-animals-breathe:2}
\omterm{def:g7:how-animals-breathe:four}{फेफड़े}: \omterm{prop:g3:sorting-living-things:groups}{स्तनधारी}, \omterm{prop:g3:sorting-living-things:groups}{पक्षी} (और \omterm{prop:g4:classifying-animals:classes}{सरीसृप}, वयस्क \omterm{prop:g3:sorting-living-things:groups}{उभयचर} भी)। \omterm{def:g7:how-animals-breathe:four}{क्लोम}: \omterm{prop:g3:sorting-living-things:groups}{मछलियाँ},
\omterm{ex:g2:animals-grow-up:frog}{टैडपोल} (और सीपियाँ, पानी के डिंभ भी)। \omterm{def:g7:how-animals-breathe:four}{श्वासनलिकाएँ}: टिड्डे, मधुमक्खियाँ
--- यानी आम तौर पर \omterm{prop:g3:sorting-living-things:groups}{कीट}। \omterm{def:g7:how-animals-breathe:four}{त्वचीय श्वसन}: केंचुआ; और \omterm{prop:g3:sorting-living-things:groups}{उभयचर}, पूरक के रूप
में।
\end{solution}

\begin{solution}{exo:g7:how-animals-breathe:3}
पानी मुँह से भीतर आता है, पंखदार \omterm{def:g7:how-animals-breathe:four}{क्लोमों} के ऊपर से धकेला जाता है ---
जहाँ विनिमय होता है --- और क्लोम-ढक्कनों के नीचे से बाहर निकल जाता है:
यानी एक-तरफ़ा धार, जिसे मुँह और ढक्कन बारी-बारी से पंप करते हैं।
\end{solution}

\begin{solution}{exo:g7:how-animals-breathe:4}
\omterm{def:g7:how-animals-breathe:four}{श्वासनलिकाओं} से: महीन हवा-नलियों का वह जाल जो शरीर के छिद्रों से चलकर हर
अंग तक शाखाएँ बनाता है --- हवा घर-घर पहुँचाई जाती है, इसलिए \omterm{prop:g4:journey-of-food:absorb}{रक्त} को
बहुत कम ढोना पड़ता है।
\end{solution}

\begin{solution}{exo:g7:how-animals-breathe:5}
पानी से बाहर उसकी पंखदार क्लोम-सतहें बैठकर आपस में चिपक जाती हैं: वह
बड़ी सतह लगभग शून्य हो जाती है, और सतह न हो तो पार करना भी नहीं ---
\omterm{def:g7:what-is-respiration:respiration}{ऑक्सीजन} भरपूर, पर उसे लेने वाली मशीन नहीं।
\end{solution}

\begin{solution}{exo:g7:how-animals-breathe:6}
\omterm{def:g7:how-animals-breathe:four}{त्वचीय श्वसन} --- और नमी वाली शर्त टूट जाती है: सूखती \omterm{def:g1:the-five-senses:sense}{त्वचा} गैसें पार
नहीं करा पाती, और केंचुआ खुली हवा में दम तोड़ देता है। जो सतह अब नम ही
न रही, उसे बड़ी और पतली होना नहीं बचा सकता।
\end{solution}

\begin{solution}{exo:g7:how-animals-breathe:7}
\omterm{ex:g2:animals-grow-up:frog}{टैडपोल}: \omterm{def:g7:how-animals-breathe:four}{क्लोम}। वयस्क: सादे \omterm{def:g7:how-animals-breathe:four}{फेफड़े} और साथ में नम \omterm{def:g1:the-five-senses:sense}{त्वचा} --- और तालाब की तह
वाली सुस्त सर्दी भर अकेली \omterm{def:g1:the-five-senses:sense}{त्वचा} ही सेवा करती है।
\end{solution}

\begin{solution}{exo:g7:how-animals-breathe:8}
उसकी मशीन \omterm{def:g7:how-animals-breathe:four}{फेफड़ा} है --- जो \omterm{def:g2:animals-grow-up:born}{दूध} पिलाने वाले थलचर पूर्वजों से विरासत में
मिला है --- इसलिए उसका जीवन चाहे कितना ही मछली जैसा हो, उसे साँस के लिए
सतह पर आना ही पड़ता है: नीचे रोक लो तो वह डूब जाती है। लिबास पानी का
है; मशीन कुल की।
\end{solution}

\begin{solution}{exo:g7:how-animals-breathe:9}
ढोया हुआ हवा-भंडार --- यानी डैनों के ढक्कनों के नीचे ग़ोताख़ोर की टंकी,
जिससे उसकी \omterm{def:g7:how-animals-breathe:four}{श्वासनलिकाएँ} खींचती हैं। उसे बुलबुला फिर से भरने के लिए
नियमित रूप से सतह पर लौटना पड़ता है।
\end{solution}

\begin{solution}{exo:g7:how-animals-breathe:10}
(क) सीपी: पानी की बाशिंदा, कोई दिखती गति नहीं, पर \omterm{def:g1:animals-around-us:covering}{खोल} के भीतर \omterm{def:g7:how-animals-breathe:four}{क्लोम}, जो
एक पंप की गई धार से धुलते हैं। (ख) टिड्डा: हवा का बाशिंदा; नन्हे छिद्रों
की क़तारें, धड़कता पेट --- यानी \omterm{def:g7:how-animals-breathe:four}{श्वासनलिकाएँ}। (ग) पनसाँप: \omterm{prop:g4:classifying-animals:classes}{सरीसृप}, और
आने-जाने वाला --- \omterm{def:g7:how-animals-breathe:four}{फेफड़े}; वह नथुने ऊपर रखकर तैरता है और साँस के लिए सतह
पर आता है।
\end{solution}

\begin{solution}{exo:g7:how-animals-breathe:11}
\omterm{def:g7:how-animals-breathe:four}{श्वासनलिकाएँ} शाखाओं वाली हवा-नलियों से पहुँचाती हैं, जो मिलीमीटरों पर
बेहतरीन और लंबी दूरियों पर बहुत धीमी हैं: ह्वेल जितने बड़े शरीर के भीतरी
हिस्से किसी भी छिद्र से बहुत दूर पड़ जाते। मशीन आकार की एक छत तय कर देती
है, और \omterm{prop:g3:sorting-living-things:groups}{कीट} उसी के नीचे जीते हैं।
\end{solution}

\begin{solution}{exo:g7:how-animals-breathe:12}
विशेष-सूची: शरीर और आस-पास की एक बड़ी, पतली, नम मिलन-सतह। \omterm{def:g7:how-animals-breathe:four}{क्लोम} इसे उन
पंखदार सतहों से पूरा करते हैं जिन्हें पानी फैलाए रखता है; \omterm{def:g7:how-animals-breathe:four}{फेफड़े} नम ओट
में तह की गई लाखों कोठरियों से; \omterm{def:g7:how-animals-breathe:four}{श्वासनलिकाएँ} अपने सबसे महीन सिरों की
जोड़ी हुई भीतरी सतह से, जो सिरों पर नम रहती है; और \omterm{def:g1:the-five-senses:sense}{त्वचा} शरीर की पूरी नम
बाहरी सतह से, जो छोटे पतले शरीर के मुक़ाबले बड़ी बनी रहती है। चार थामने
वाले, और थामी हुई चीज़ एक।
\end{solution}

\begin{solution}{pb:g7:how-animals-breathe:1}
\textbf{1.} ट्राउट के बच्चे और काँटा-मछलियाँ: \omterm{def:g7:how-animals-breathe:four}{क्लोम}। \omterm{ex:g2:animals-grow-up:frog}{टैडपोल}: \omterm{def:g7:how-animals-breathe:four}{क्लोम}
(\omterm{def:g1:my-body:parts}{टाँगों} वाले पुनर्निर्माण के बीच में)। वयस्क मेंढक: \omterm{def:g7:how-animals-breathe:four}{फेफड़े} और साथ में
\omterm{def:g1:the-five-senses:sense}{त्वचा}। तालाब के घोंघे: \omterm{def:g7:how-animals-breathe:four}{फेफड़े} --- सतह पर आने-जाने वाले। गोताख़ोर भृंग:
\omterm{def:g7:how-animals-breathe:four}{श्वासनलिकाएँ}, जो उसके ढोए बुलबुले से भरती हैं। मेफ़्लाई के डिंभ: \omterm{def:g7:how-animals-breathe:four}{क्लोम}
(वही पंखदार गुच्छे)। केंचुए: \omterm{def:g1:the-five-senses:sense}{त्वचा}।

\textbf{2.} \omterm{def:g7:how-animals-breathe:four}{क्लोम}। उनका हिलना क्लोम-सतह के सामने का पानी नया करता रहता
है --- यानी ख़ुद बनाई हुई धार, वहाँ जहाँ ठहरा तालाब कोई नहीं देता, ताकि
पार करने की जगह पर \omterm{def:g7:what-is-respiration:respiration}{ऑक्सीजन} वाला ताज़ा पानी बना रहे।

\textbf{3.} तालाब के घोंघे, गोताख़ोर भृंग, वयस्क मेंढक (और आने वाला कोई
पनसाँप भी): यानी \omterm{def:g7:how-animals-breathe:four}{फेफड़े} और बुलबुले वाले सब --- उनकी मशीनें \omterm{def:g7:what-is-respiration:respiration}{ऑक्सीजन} हवा
से लेती हैं, इसलिए सतह ही उनका भरने का अड्डा है।

\textbf{4.} \omterm{def:g2:animals-grow-up:metamorphosis}{कायांतरण} का साँस वाला पुनर्निर्माण: बिना \omterm{def:g1:my-body:parts}{टाँगों} वाले \omterm{ex:g2:animals-grow-up:frog}{टैडपोल}
की सेवा करते \omterm{def:g7:how-animals-breathe:four}{क्लोमों} की जगह वे \omterm{def:g7:how-animals-breathe:four}{फेफड़े} (और साँस लेती \omterm{def:g1:the-five-senses:sense}{त्वचा}) ली जा रही हैं
जिनकी ज़रूरत किनारे पर जाने वाले नन्हे मेंढक को पड़ेगी।

\textbf{5.} पहले तकलीफ़ \omterm{def:g7:how-animals-breathe:four}{क्लोम} वालों को होती है: उनकी पूरी आपूर्ति पानी
का घटता हुआ घुला भंडार ही है। और आने-जाने वाले --- घोंघा, भृंग, मेंढक
--- ऊपर की असीम हवा पर टिके हैं और बस ज़्यादा बार चक्कर लगा लेते हैं।

\textbf{6.} सबसे ऊपरी मिलीमीटर हवा को छूता है: ऊपर से \omterm{def:g7:what-is-respiration:respiration}{ऑक्सीजन} उसमें घुलती
रहती है, इसलिए दम घोंटते तालाब में परत वाली परत ही सबसे कम ग़रीब पानी
है। काँटा-मछलियाँ आख़िरी कुएँ पर क़तार लगाए हैं।

\textbf{7.} \omterm{def:g7:how-animals-breathe:four}{क्लोम} जाम होते हैं: महीन कीचड़ पंखदार सतहों पर बैठ जाती है
और पार करने का रास्ता रोक देती है। उम्मीद करो कि \omterm{def:g7:how-animals-breathe:four}{क्लोम} से साँस लेने वाले
उस बादल से भागेंगे, और ढक्कनों का पंप करना मेहनत करने लगेगा --- गाद हटाने
के लिए खाँसी जैसे झटके।

\textbf{8.} \omterm{def:g7:how-animals-breathe:four}{त्वचीय श्वसन}, और नम-पर-हवा-रहित वाली परिस्थिति: पानी से भरी
बिलों में इतनी कम \omterm{def:g7:what-is-respiration:respiration}{ऑक्सीजन} बचती है कि केंचुए साँस के लिए सतह पर आ जाते हैं
--- यानी नीचे दम घुटने के बदले ऊपर सूखने का ख़तरा मोल लेते हुए।

\textbf{9.} तालिका जैसा कहा गया है (हर \omterm{def:g1:animals-around-us:animal}{जंतु} की मशीन और माध्यम;
आने-जाने वाले: घोंघा, भृंग, मेंढक)। हर स्तंभ के \omterm{def:g1:my-body:parts}{सिर} पर:
बड़ी-पतली-नम विनिमय-सतह --- यानी वही एक विशेष-सूची जिसे चारों मशीनें
थामे रहती हैं।

\textbf{10.} मेंढक: इंतज़ाम है --- कीचड़ में सुस्त पड़े, अपनी छोटी सर्दी
वाली ज़रूरत \omterm{def:g1:the-five-senses:sense}{त्वचा} से पूरी करते हुए। घोंघे और भृंग: मुश्किल में ---
उनकी सतह बंद है; वे घटी हुई ज़रूरतों के साथ सुप्त रहते हैं या फँसे हुए
बुलबुलों के सहारे। \omterm{prop:g3:sorting-living-things:groups}{मछलियाँ}: तब तक इंतज़ाम है जब तक पानी का भंडार चले
--- ठंड हर इंजन की आवाज़ धीमी कर देती है
(\cref{ex:g7:what-is-respiration:demand}), और \omterm{def:g7:how-animals-breathe:four}{क्लोम} बर्फ़ के नीचे उस
घटी हुई माँग की सेवा कर लेते हैं।

\textbf{11.} मकड़ी की रेशमी घंटी एक जमी हुई हवा-टंकी है, जबकि भृंग का
बुलबुला ढोई हुई टंकी: दोनों सतह की हवा पानी के नीचे जमा रखते हैं, और
दोनों ग़लत माध्यम में एक हवा-मशीन (मकड़ी का पुस्तक-फुप्फुस और
श्वासनलिका वाला साज़-सामान) की सेवा जारी रखते हैं।

\textbf{12.} क्योंकि तालाब \omterm{def:g7:what-is-respiration:respiration}{ऑक्सीजन} दो माध्यमों में देता है --- नीचे विरल
और घुली हुई, ऊपर भरपूर --- और हर मशीन एक माध्यम को, जीवन के एक ही आकार
पर बटोरती है: चारों के काम पर रहने से ही कीचड़ से लेकर सतह की परत तक
का हर कोना बसा रह पाता है।
\end{solution}
